```latex
\documentclass[11pt]{article}
\usepackage{amsmath,amssymb,amsthm}
\usepackage[colorlinks=true,urlcolor=blue]{hyperref}

\newtheorem{theorem}{Theorem}
\newtheorem{lemma}{Lemma}
\newtheorem{proposition}{Proposition}
\newtheorem{corollary}{Corollary}
\theoremstyle{definition}
\newtheorem{definition}{Definition}

\newcommand{\Uq}{U_q^{\mathrm{ext}}}
\newcommand{\slie}{\mathfrak{sl}}
\newcommand{\fb}{\mathfrak{b}}
\newcommand{\F}{\mathcal{F}}
\newcommand{\cS}{\mathcal{S}}
\newcommand{\bipr}{\bowtie}
\newcommand{\KK}{\hat K}
\newcommand{\eps}{\varepsilon}

\title{The quantum double of Borel subalgebras and the bigon presentation of $\Uq(n)$}
\author{}
\date{}

\begin{document}

\maketitle

Throughout, $k$ is a field of characteristic $0$, $q\in k^\times$ is not a root of unity,
and all algebras are over $k(q)$. We write
\[
[n]_q=\frac{q^n-q^{-n}}{q-q^{-1}}.
\]
We use the simply-connected quantized enveloping algebra $\Uq(\slie(n))$ as in
Klimyk--Schm\"udgen. Its Cartan part has group-like generators
$\KK_1^{\pm1},\dots,\KK_n^{\pm1}$ with
\[
\KK_1\cdots \KK_n=1.
\]
For $1\le i<n$, put
\[
K_i:=\KK_i \KK_{i+1}^{-1}.
\]
The Chevalley generators $E_i,F_i$ satisfy
\begin{equation}
  \KK_j E_i=q^{\delta_{ji}-\delta_{j,i+1}}E_i \KK_j,\qquad
  \KK_j F_i=q^{-(\delta_{ji}-\delta_{j,i+1})}F_i \KK_j,
\end{equation}
\begin{equation}
  E_iF_j-F_jE_i=\delta_{ij}\frac{K_i-K_i^{-1}}{q-q^{-1}},
\end{equation}
together with the usual $q$-Serre relations.

Let $\Uq(\fb_+)$ and $\Uq(\fb_-)$ be the upper and lower Borel subalgebras:
\[
\Uq(\fb_+)=\bigl\langle \KK_i^{\pm1}, E_i : 1\le i<n\bigr\rangle,\qquad
\Uq(\fb_-)=\bigl\langle \KK_i^{\pm1}, F_i : 1\le i<n\bigr\rangle.
\]
Their Hopf structures are the restrictions of the Hopf structure of $\Uq(\slie(n))$:
\begin{align*}
\Delta(\KK_i)&=\KK_i\otimes \KK_i,\\
\Delta(E_i)&=E_i\otimes 1+K_i\otimes E_i,\qquad S(E_i)=-K_i^{-1}E_i,\\
\Delta(F_i)&=F_i\otimes K_i^{-1}+1\otimes F_i,\qquad S(F_i)=-F_iK_i.
\end{align*}

\section{The Drinfeld--Tanisaki pairing}

\begin{definition}
The standard Drinfeld--Tanisaki Hopf pairing
\[
\phi:\Uq(\fb_+)\times \Uq(\fb_-)\longrightarrow k(q)
\]
is the unique Hopf pairing determined by
\begin{equation}
\label{eq:pairing}
\phi(\KK_i,\KK_j)=q^{\frac1n-\delta_{ij}},\qquad
\phi(E_i,F_j)=\frac{\delta_{ij}}{q^{-1}-q},\qquad
\phi(\KK_i,F_j)=\phi(E_j,\KK_i)=0.
\end{equation}
\end{definition}

That this is indeed a well-defined nondegenerate Hopf pairing is the standard
Drinfeld--Tanisaki pairing; see for example
\href{https://staff.fnwi.uva.nl/j.v.stokman/Quantumsl2.pdf}{staff.fnwi.uva.nl}
for the $\slie(2)$ case and the general construction in
\href{https://www.cambridge.org/core/services/aop-cambridge-core/content/view/886B0B4D8E3D242A51A0380A486351EA/S0013091504000689a.pdf/presentation-by-borel-subalgebras-and-chevalley-generators-for-quantum-enveloping-algebras.pdf}{cambridge.org}.
For later use, note that
\[
\phi(K_i,K_j)=q^{-a_{ij}},
\]
where $(a_{ij})$ is the Cartan matrix of $\slie(n)$.

\section{The generalized quantum double}

Let $A,B$ be Hopf algebras and let $\phi:A\times B\to k$ be a Hopf pairing.
The generalized quantum double $D_\phi(A,B)$ is the algebra generated by the
two vector-space copies $A\otimes 1$ and $1\otimes B$, with cross relation
\begin{equation}
\label{eq:cross-general}
(1\otimes b)(a\otimes 1)
=
\sum_{(a),(b)}
\phi\bigl(a_{(1)},S_B(b_{(1)})\bigr)
\phi\bigl(a_{(3)},b_{(3)}\bigr)
(a_{(2)}\otimes 1)(1\otimes b_{(2)}),
\end{equation}
where $\Delta_A^2(a)=\sum a_{(1)}\otimes a_{(2)}\otimes a_{(3)}$ and similarly
for $b$. This is the algebraic form of the Drinfeld double construction; see
\href{https://hal.science/hal-02159493v1/file/CQG.pdf}{hal.science}
and
\href{https://staff.fnwi.uva.nl/j.v.stokman/quantumdouble2.pdf}{staff.fnwi.uva.nl}.

We work with
\[
D:=D_\phi\bigl(\Uq(\fb_+),\Uq(\fb_-)\bigr).
\]
Up to the flip of tensor factors, this is the algebra denoted
$\Uq(\fb_-)\bipr\Uq(\fb_+)$ in the statement.

The following computation gives the defining cross relations in $D$.

\begin{lemma}
\label{lem:cross}
In $D$, for $1\le i,j<n$ and $1\le r,s\le n$, we have
\begin{align}
(\KK_r\otimes 1)(1\otimes \KK_s)
&=
(1\otimes \KK_s)(\KK_r\otimes 1),
\label{eq:c0}
\\
(1\otimes \KK_s)(E_i\otimes 1)
&=
q^{\delta_{si}-\delta_{s,i+1}}
(E_i\otimes 1)(1\otimes \KK_s),
\label{eq:c1}
\\
(\KK_r\otimes 1)(1\otimes F_j)
&=
q^{\delta_{r,j+1}-\delta_{rj}}
(1\otimes F_j)(\KK_r\otimes 1),
\label{eq:c2}
\\
(E_i\otimes 1)(1\otimes F_j)
-
(1\otimes F_j)(E_i\otimes 1)
&=
\delta_{ij}\,
\frac{K_i\otimes 1-1\otimes K_j^{-1}}{q-q^{-1}}.
\label{eq:c3}
\end{align}
\end{lemma}

\begin{proof}
The identities follow from \eqref{eq:cross-general}.

For \eqref{eq:c0}, both $a=\KK_r$ and $b=\KK_s$ are group-like, so
$\Delta_A^2(a)=a\otimes a\otimes a$ and
$\Delta_B^2(b)=b\otimes b\otimes b$. Thus
\[
(1\otimes \KK_s)(\KK_r\otimes 1)
=
\phi(\KK_r,\KK_s^{-1})\phi(\KK_r,\KK_s)
(\KK_r\otimes 1)(1\otimes \KK_s)
=
(\KK_r\otimes 1)(1\otimes \KK_s),
\]
because the two scalar factors are inverse to each other.

For \eqref{eq:c1}, take $a=E_i$, $b=\KK_s$. Since $b$ is group-like and
\[
\Delta_A^2(E_i)
=
E_i\otimes 1\otimes 1
+
K_i\otimes E_i\otimes 1
+
K_i\otimes K_i\otimes E_i,
\]
the only nonvanishing term in \eqref{eq:cross-general} is the term containing
$\phi(K_i,\KK_s^{-1})$. Therefore
\[
(1\otimes \KK_s)(E_i\otimes 1)
=
\phi(K_i,\KK_s^{-1})(E_i\otimes 1)(1\otimes \KK_s).
\]
Using \eqref{eq:pairing},
\[
\phi(K_i,\KK_s^{-1})
=
q^{\delta_{si}-\delta_{s,i+1}},
\]
which proves \eqref{eq:c1}.

For \eqref{eq:c2}, take $a=\KK_r$, $b=F_j$. Now $a$ is group-like and
\[
\Delta_B^2(F_j)
=
F_j\otimes K_j^{-1}\otimes K_j^{-1}
+
1\otimes F_j\otimes K_j^{-1}
+
1\otimes 1\otimes F_j.
\]
Since $\phi(\KK_r,F_j)=0$ and $\phi(\KK_r,S(F_j))=0$, the only contributing term is
\[
(1\otimes F_j)(\KK_r\otimes 1)
=
\phi(\KK_r,K_j^{-1})(\KK_r\otimes 1)(1\otimes F_j).
\]
Again by \eqref{eq:pairing},
\[
\phi(\KK_r,K_j^{-1})=q^{\delta_{rj}-\delta_{r,j+1}}.
\]
Rearranging gives \eqref{eq:c2}.

Finally, for \eqref{eq:c3}, take $a=E_i$, $b=F_j$. A direct expansion of
\eqref{eq:cross-general}, using the above formulas for $\Delta_A^2(E_i)$ and
$\Delta_B^2(F_j)$, gives
\[
(1\otimes F_j)(E_i\otimes 1)
=
(E_i\otimes 1)(1\otimes F_j)
+
\frac{\delta_{ij}}{q^{-1}-q}
\bigl(K_i\otimes 1-1\otimes K_j^{-1}\bigr).
\]
Multiplying by $-1$ and using $q^{-1}-q=-(q-q^{-1})$ yields \eqref{eq:c3}.
\end{proof}

\section{Root vectors and the PBW basis}

For $1\le i<n$, put
\[
E_{i,i+1}=E_i,\qquad F_{i,i+1}=F_i.
\]
For $j>i+1$, define recursively
\begin{equation}
\label{eq:upper-root}
E_{ij}=E_{i,j-1}E_{j-1,j}-q^{-1}E_{j-1,j}E_{i,j-1},
\end{equation}
and
\begin{equation}
\label{eq:lower-root}
F_{ij}=F_{i,j-1}F_{j-1,j}-q F_{j-1,j}F_{i,j-1}.
\end{equation}
The elements $E_{ij}$ are the standard positive root vectors for the root
$\eps_i-\eps_j$, and the $F_{ij}$ are the corresponding negative root vectors.
The signs are chosen so that the two families are dual under the pairing $\phi$;
for $\slie(3)$ this is the familiar
\[
E_{13}=E_1E_2-q^{-1}E_2E_1,\qquad
F_{13}=F_1F_2-qF_2F_1.
\]

By the PBW theorem for quantized enveloping algebras, the set
\[
\Bigl\{\prod_{i<j}E_{ij}^{m_{ij}}\prod_{r=1}^n \KK_r^{a_r}
: m_{ij}\in\mathbb Z_{\ge0},\ a_r\in\mathbb Z\Bigr\}
\]
is a basis of $\Uq(\fb_+)$, and similarly the set
\[
\Bigl\{\prod_{r=1}^n \KK_r^{b_r}\prod_{i<j}F_{ij}^{n_{ij}}
: n_{ij}\in\mathbb Z_{\ge0},\ b_r\in\mathbb Z\Bigr\}
\]
is a basis of $\Uq(\fb_-)$, with the products taken in any fixed convex order.
See
\href{https://www.cambridge.org/core/services/aop-cambridge-core/content/view/886B0B4D8E3D242A51A0380A486351EA/S0013091504000689a.pdf/presentation-by-borel-subalgebras-and-chevalley-generators-for-quantum-enveloping-algebras.pdf}{cambridge.org}
for the general Borel-subalgebra presentation.

Since $D=\Uq(\fb_+)\otimes\Uq(\fb_-)$ as a vector space, we obtain the following
PBW basis of the double.

\begin{proposition}
\label{prop:pbw}
The double $D$ has a $k(q)$-basis consisting of the ordered monomials
\begin{equation}
\label{eq:pbw-double}
\prod_{i<j}(E_{ij}\otimes 1)^{m_{ij}}
\prod_{r=1}^n(\KK_r\otimes 1)^{a_r}
\prod_{r=1}^n(1\otimes \KK_r^{-1})^{b_r}
\prod_{i<j}(1\otimes F_{ij})^{n_{ij}},
\end{equation}
where $m_{ij},n_{ij}\in\mathbb Z_{\ge0}$ and $a_r,b_r\in\mathbb Z$.
\end{proposition}

\section{The algebra $\cS_n(\F^*)$}

The diagrammatic algebra $\cS_n(\F^*)$ is generated by the horizontal arcs
$b_{ij},c_{ij}$, $1\le i\le j\le n$, directed to the right. Its defining bigon
relations are equivalent to the following algebraic presentation.

\begin{definition}
\label{def:S}
The algebra $\cS_n(\F^*)$ is generated by $b_{ij}$ and $c_{ij}$ for
$1\le i\le j\le n$, with $b_{ii}$ and $c_{ii}$ invertible, subject to
the following relations.

\smallskip
\noindent\textbf{(i) Diagonal relations.}
\[
b_{ii}b_{jj}=b_{jj}b_{ii},\qquad
c_{ii}c_{jj}=c_{jj}c_{ii},\qquad
b_{ii}c_{jj}=c_{jj}b_{ii},
\]
\[
b_{11}\cdots b_{nn}=1,\qquad c_{11}\cdots c_{nn}=1.
\]

\smallskip
\noindent\textbf{(ii) Root-vector relations.}
For $i<j<k$,
\begin{equation}
\label{eq:S-upper-root}
b_{ik}=b_{ij}b_{jk}-q^{-1}b_{jk}b_{ij},
\end{equation}
and
\begin{equation}
\label{eq:S-lower-root}
c_{ik}=c_{ij}c_{jk}-q c_{jk}c_{ij}.
\end{equation}

\smallskip
\noindent\textbf{(iii) Internal Borel relations.}
For $1\le r\le n$ and $i<j$,
\begin{equation}
\label{eq:S-upper-diag}
b_{rr}b_{ij}=q^{\delta_{ri}-\delta_{rj}}b_{ij}b_{rr},
\end{equation}
\begin{equation}
\label{eq:S-lower-diag}
c_{rr}c_{ij}=q^{\delta_{ri}-\delta_{rj}}c_{ij}c_{rr}.
\end{equation}
The simple generators $b_{i,i+1}$ satisfy the usual $q$-Serre relations of
$\Uq(\fb_+)$, and the $c_{i,i+1}$ satisfy the $q$-Serre relations of
$\Uq(\fb_-)$.

\smallskip
\noindent\textbf{(iv) Cross relations.}
For $1\le r\le n$ and $i<j$,
\begin{equation}
\label{eq:S-cross-cartan}
c_{rr}b_{ij}=q^{\delta_{rj}-\delta_{ri}}b_{ij}c_{rr},
\end{equation}
\begin{equation}
\label{eq:S-cross-cartan2}
b_{rr}c_{ij}=q^{\delta_{rj}-\delta_{ri}}c_{ij}b_{rr}.
\end{equation}
For $1\le i,j<n$,
\begin{equation}
\label{eq:S-simple-comm}
b_{i,i+1}c_{j,j+1}-c_{j,j+1}b_{i,i+1}
=
\delta_{ij}\,
\frac{b_{ii}b_{i+1,i+1}^{-1}-c_{ii}c_{i+1,i+1}^{-1}}
     {q-q^{-1}}.
\end{equation}
\end{definition}

\section{The explicit isomorphism}

\begin{theorem}
\label{thm:isom}
There is an algebra isomorphism
\[
\Phi:\cS_n(\F^*)
\overset{\sim}{\longrightarrow}
D_\phi\bigl(\Uq(\fb_+),\Uq(\fb_-)\bigr)
\cong
\Uq(\fb_-)\bipr\Uq(\fb_+),
\]
determined by
\begin{align*}
b_{ii}&\longmapsto \KK_i\otimes 1,
&
c_{ii}&\longmapsto 1\otimes \KK_i^{-1},
\\
b_{ij}&\longmapsto E_{ij}\otimes 1,
&
c_{ij}&\longmapsto 1\otimes F_{ij},
\qquad i<j.
\end{align*}
\end{theorem}

\begin{proof}
We first verify that $\Phi$ is well defined.

The relations in Definition~\ref{def:S}(i)--(iii) are exactly the defining
relations of the two Borel subalgebras, with $\KK_i^{-1}$ used as the diagonal
generator in $\Uq(\fb_-)$. Hence they are preserved by $\Phi$.

The cross relations \eqref{eq:S-cross-cartan} and \eqref{eq:S-cross-cartan2}
follow from Lemma~\ref{lem:cross}, \eqref{eq:c1} and \eqref{eq:c2}, after
passing from simple root vectors to arbitrary root vectors using the formulas
\eqref{eq:upper-root} and \eqref{eq:lower-root}. Finally, the simple commutator
\eqref{eq:S-simple-comm} is exactly Lemma~\ref{lem:cross}, \eqref{eq:c3},
because
\[
K_i\otimes 1
=
b_{ii}b_{i+1,i+1}^{-1},
\qquad
1\otimes K_j^{-1}
=
c_{jj}c_{j+1,j+1}^{-1}.
\]
Thus $\Phi$ is a well-defined algebra homomorphism.

It remains to show that $\Phi$ is bijective. Surjectivity is clear, since the
images of the diagonal arcs give the two Cartan copies, and the images of
$b_{i,i+1}$ and $c_{i,i+1}$ give the Chevalley generators $E_i\otimes 1$ and
$1\otimes F_i$.

For injectivity, use the defining relations of $\cS_n(\F^*)$ to rewrite any
word in the generators as a linear combination of ordered monomials
\[
\prod_{i<j}b_{ij}^{m_{ij}}
\prod_{r=1}^n b_{rr}^{a_r}
\prod_{r=1}^n c_{rr}^{b_r}
\prod_{i<j}c_{ij}^{n_{ij}}.
\]
This is the usual PBW reduction argument for Drinfeld doubles: the internal
relations \eqref{eq:S-upper-root}--\eqref{eq:S-lower-diag} give PBW bases inside
the two Borel parts, while the cross relations
\eqref{eq:S-cross-cartan}--\eqref{eq:S-simple-comm} allow one to move all
$c$-generators to the right of all $b$-generators at the expense of producing
diagonal terms.

Under $\Phi$, these ordered monomials are sent exactly to the PBW basis
\eqref{eq:pbw-double} of Proposition~\ref{prop:pbw}. Hence the images are
linearly independent, and therefore $\Phi$ is injective.
\end{proof}

\section{The quotient by $b_{ii}c_{ii}-1$}

Let
\[
I:=\bigl\langle b_{ii}c_{ii}-1 : 1\le i\le n\bigr\rangle
\subset \cS_n(\F^*).
\]
Under the isomorphism $\Phi$, this ideal corresponds to the two-sided ideal
\[
J:=\bigl\langle \KK_i\otimes \KK_i^{-1}-1 : 1\le i\le n\bigr\rangle
\subset D.
\]

\begin{theorem}
\label{thm:quotient}
There is an algebra isomorphism
\[
\cS_n(\F^*)/I
\cong
\Uq(\slie(n)).
\]
\end{theorem}

\begin{proof}
It suffices to prove that
\[
D/J\cong \Uq(\slie(n)).
\]
This is the $\slie(n)$ generalization of the $\slie(2)$ computation in
\href{https://staff.fnwi.uva.nl/j.v.stokman/quantumdouble2.pdf}{staff.fnwi.uva.nl},
Corollary 2.12.

Define an algebra homomorphism
\[
\Theta:D\longrightarrow \Uq(\slie(n))
\]
by
\[
E_i\otimes 1\mapsto E_i,\qquad
1\otimes F_i\mapsto F_i,\qquad
\KK_r\otimes 1\mapsto \KK_r,\qquad
1\otimes \KK_r\mapsto \KK_r.
\]
The relations of $D$ are exactly the relations of $\Uq(\slie(n))$ after this
identification: the cross commutator \eqref{eq:c3} becomes
\[
E_iF_j-F_jE_i=\delta_{ij}\frac{K_i-K_i^{-1}}{q-q^{-1}},
\]
and the cross Cartan relations \eqref{eq:c1}, \eqref{eq:c2} become the
standard Cartan--Chevalley relations. Thus $\Theta$ is well defined.

The kernel of $\Theta$ contains $J$, because
\[
\Theta\bigl(\KK_i\otimes \KK_i^{-1}-1\bigr)
=
\KK_i\KK_i^{-1}-1=0.
\]
Hence $\Theta$ descends to an algebra homomorphism
\[
\overline\Theta:D/J\longrightarrow \Uq(\slie(n)).
\]

Conversely, define
\[
\Psi:\Uq(\slie(n))\longrightarrow D/J
\]
by
\[
\KK_r\mapsto [\KK_r\otimes 1]=[1\otimes \KK_r],\qquad
E_i\mapsto [E_i\otimes 1],\qquad
F_i\mapsto [1\otimes F_i].
\]
The equality $[\KK_r\otimes 1]=[1\otimes \KK_r]$ holds in $D/J$ precisely because
$J$ contains $\KK_r\otimes \KK_r^{-1}-1$. The defining relations of
$\Uq(\slie(n))$ are satisfied in $D/J$ by Lemma~\ref{lem:cross} and the Borel
subalgebra relations. Thus $\Psi$ is well defined.

On generators, $\overline\Theta\circ\Psi$ and
$\Psi\circ\overline\Theta$ are the identity maps. Therefore
\[
D/J\cong \Uq(\slie(n)).
\]
Combining this with Theorem~\ref{thm:isom} gives
\[
\cS_n(\F^*)/I\cong D/J\cong \Uq(\slie(n)).
\]
\end{proof}

\begin{corollary}
The algebra $\cS_n(\F^*)$ is a flat deformation of the algebra obtained from
the two Borel subalgebras of $\Uq(\slie(n))$ by identifying the two Cartan
copies. In particular, its quotient by the diagonal arc relations
$b_{ii}c_{ii}=1$ recovers the simply-connected quantized enveloping algebra
$\Uq(\slie(n))$.
\end{corollary}

\end{document}
```